\chapter{Fonds indiciels cotés (ETF)}\label{ch:m1:exchange-traded-funds}

Les parts d'un fonds s'échangent à 50.10~\$ en bourse alors que les cinq cents
actions qu'il renferme valent 50.00~\$ l'action. En une seconde ou deux, une
société de trading vend les parts du fonds, achète les cinq cents actions et, le soir même, remet ces
actions au fonds, en échange de parts nouvellement émises qui soldent
sa position vendeuse. Elle a gagné quelques centimes par part~; le fonds a grandi~; le cours
est de nouveau aligné. Personne au sein du fonds n'a rien décidé. Fin août 2026
quelque 24~billions de dollars étaient investis dans des véhicules tenus honnêtes par
ce seul mécanisme. Ce chapitre explique ce mécanisme et ce qui lui arrive
lorsque les actions qu'il contient sont fermées ou ne peuvent pas être évaluées, et les fonds
qui promettent trois fois l'indice et livrent tout autre chose.

\section{Un fonds à deux marchés}

\begin{definition}[Fonds indiciel coté (ETF)]\label{def:m1:exchange-traded-funds:etf}
Un \emph{fonds indiciel coté}\index{fonds indiciel coté} (exchange-traded fund, ETF) est un
fonds d'investissement ouvert dont les parts sont cotées et se négocient en bourse
tout au long de la journée, comme une action quelconque, et peuvent en outre être
créées ou rachetées auprès du fonds lui-même, par blocs importants, par des intermédiaires désignés.
\end{definition}

\begin{definition}[Participant autorisé, unité de création, panier de création]\label{def:m1:exchange-traded-funds:ap}
Un \emph{participant autorisé}\index{participant autorisé} (authorised participant, AP) est un
\omterm{def:m1:the-sell-side:sell-side}{broker-dealer} ayant conclu avec un ETF un contrat lui donnant le droit de créer et
racheter ses parts. Il le fait par \emph{unités de création}\index{unité de création},
blocs d'un nombre fixe de parts d'ETF, en livrant au fonds (ou
en recevant de lui) le \emph{panier de création}\index{panier de création}~: la
liste des titres et des liquidités publiée chaque jour par le fonds, et
contre laquelle une unité de création est échangée.
\end{definition}

\begin{definition}[Transfert en nature]\label{def:m1:exchange-traded-funds:inkind}
Une création ou un rachat est \emph{en nature}\index{transfert en nature} lorsque
l'AP et le fonds échangent des parts d'ETF contre les titres eux-mêmes, et non
contre des liquidités. Le fonds n'opère alors jamais sur le marché~: le coût de
la transformation de liquidités en portefeuille, et inversement, est supporté par les
investisseurs qui entrent et sortent, et non par ceux qui restent.
\end{definition}

\begin{omfigure}
\begin{tikzpicture}[font=\small,
  box/.style={draw=omDef,rounded corners=2pt,minimum width=2.9cm,minimum height=1.0cm,align=center,fill=omDef!4},
  ext/.style={draw=omExo,rounded corners=2pt,minimum width=2.9cm,minimum height=1.0cm,align=center,fill=omExo!6},
  flow/.style={-{Stealth[length=2mm]},thick}]
\node[ext] (inv) at (0,1.4) {Investisseurs};
\node[box] (ex) at (5.2,1.4) {Bourse\\(marché secondaire)};
\node[box] (ap) at (10.4,1.4) {Participants autorisés\\et teneurs de marché};
\node[box] (fund) at (10.4,-1.6) {Le fonds\\(marché primaire)};
\node[ext] (und) at (5.2,-1.6) {Marché des\\titres sous-jacents};
\draw[flow,omDef] ([yshift=5pt]inv.east) -- node[above,font=\footnotesize]{transaction} ([yshift=5pt]ex.west);
\draw[flow,omDef] ([yshift=5pt]ap.west) -- node[above,font=\footnotesize]{cotation} ([yshift=5pt]ex.east);
\draw[flow,omThm] ([xshift=-7pt]ap.south) -- node[left,font=\footnotesize,align=right]{panier entrant,\\parts sortantes} ([xshift=-7pt]fund.north);
\draw[flow,omThm] ([xshift=7pt]fund.north) -- node[right,font=\footnotesize,align=left]{ou l'inverse} ([xshift=7pt]ap.south);
\draw[flow,omDef] ([yshift=-8pt]ap.west) -- ++(-1.0,0) |- node[pos=0.25,left,font=\footnotesize,align=right]{achète ou vend\\le panier} (und.east);
\end{tikzpicture}

\omcaption{Les deux marchés d'un ETF. Les investisseurs ne rencontrent que la bourse. Le
nombre de parts d'ETF en circulation ne varie que du côté droit, par
\omterm{def:m1:exchange-traded-funds:ap}{unités de création}, à la \omterm{def:m1:exchange-traded-funds:nav}{valeur liquidative} du jour, contre le
panier.}\label{fig:m1:exchange-traded-funds:two}
\end{omfigure}

\begin{dated}{2026-09}{Taille, et la règle qui a normalisé la structure}\label{dat:m1:exchange-traded-funds:size}
Les encours des \omterm{def:m1:exchange-traded-funds:etf}{fonds indiciels cotés} dans le monde ont atteint 24.03~billions de dollars à la fin
d'août 2026, un record, selon le décompte d'un cabinet de conseil. Aux États-Unis, la
Rule 6c-11 de l'Investment Company Act, adoptée le 25 septembre 2019, permet à un ETF
transparent de fonctionner sans ordonnance d'exemption individuelle, à condition qu'il publie
son portefeuille chaque jour, et autorise des \emph{paniers personnalisés} qui s'écartent
d'une tranche au prorata du portefeuille. Les \omterm{def:m1:exchange-traded-funds:leveraged}{ETF à effet de levier} et les ETF inverses ne peuvent pas invoquer cette règle.
\end{dated}

\section{L'arbitrage qui maintient le cours}

\begin{definition}[Valeur liquidative, valeur liquidative indicative, prime]\label{def:m1:exchange-traded-funds:nav}
La \emph{valeur liquidative}\index{valeur liquidative} (net asset value, NAV) d'une part de fonds est
la valeur des actifs du fonds diminuée de ses passifs, divisée par le nombre
de parts, calculée une fois par jour aux cours de clôture officiels. La
\emph{valeur liquidative indicative}\index{valeur liquidative indicative} (indicative NAV, iNAV) est une estimation de la même
grandeur recalculée au cours de la journée à partir des cours actuels du panier. La
\emph{prime}\index{prime} d'un ETF est $P/\mathrm{NAV} - 1$, avec $P$ son
cours de marché~; une prime négative est une \emph{décote}.
\end{definition}

\begin{proposition}[La bande d'arbitrage]\label{prop:m1:exchange-traded-funds:band}
Soient $c_B$ et $c_E$ les coûts de négociation du panier et de l'ETF (la moitié de la
fourchette plus l'impact, en fractions de la valeur), $\phi$ les frais de création ou de rachat
par unité de valeur, et $f$ le coût de financement des deux jambes jusqu'à la
fin du règlement de la création. Un \omterm{def:m1:exchange-traded-funds:ap}{participant autorisé} tire profit de la création lorsque
la \omterm{def:m1:exchange-traded-funds:nav}{prime} dépasse $b = c_B + c_E + \phi + f$, et de même pour le rachat lorsque la
décote dépasse $b$. Dans un marché où les AP se font concurrence, la \omterm{def:m1:exchange-traded-funds:nav}{prime}
reste donc dans l'intervalle $[-b, +b]$.
\end{proposition}
\begin{proof}
Pour une \omterm{def:m1:exchange-traded-funds:nav}{prime} $\pi > 0$, l'AP vend des parts d'ETF à $P(1 - c_E)$ et achète le
panier à $\mathrm{NAV}(1 + c_B)$, acquitte les frais et le financement, et
livre le panier en échange de parts d'ETF à la \omterm{def:m1:exchange-traded-funds:nav}{valeur liquidative}. Son profit par unité de valeur est
$\pi - c_E - c_B - \phi - f$ au premier ordre, positif si et seulement si $\pi > b$. La
opération de rachat est symétrique.
\end{proof}

\begin{example}[Sept points de base]\label{ex:m1:exchange-traded-funds:band}
Pour un ETF d'actions à grande capitalisation, prenons $c_B = 4$, $c_E = 1.5$, $\phi = 1$
et $f = 0.5$ points de base~: $b = 7$ points de base, soit trois centimes et demi sur une
action de 50~\$. Pour un fonds d'actions de petites capitalisations ou étrangères, $c_B$ seul peut atteindre trente
points de base~; pour un fonds d'obligations d'entreprises, davantage. La bande n'est pas un défaut
de l'ETF~: c'est le coût de négociation du sous-jacent, que l'investisseur en ETF
évite tant que le cours reste à l'intérieur
(la \cref{fig:m1:exchange-traded-funds:premium}).
\end{example}

\begin{omfigure}
\begin{tikzpicture}
\begin{axis}[width=11cm,height=4.8cm,
  xlabel={temps (pas arbitraires)},ylabel={prime (points de base)},
  tick label style={font=\small},label style={font=\small},
  xmin=0,xmax=120,ymin=-10,ymax=10,grid=major,grid style={gray!25},table/col sep=comma]
\addplot[omDef,thick,restrict x to domain=0:120] table[x=t,y=premium_bp]{figdata/markets-1/14-exchange-traded-funds/premium.csv};
\addplot[omThm,dashed,thick] coordinates {(0,7) (120,7)};
\addplot[omThm,dashed,thick] coordinates {(0,-7) (120,-7)};
\end{axis}
\end{tikzpicture}

\omcaption{Une \omterm{def:m1:exchange-traded-funds:nav}{prime} simulée. Le flux d'ordres la fait osciller~; chaque fois qu'elle
atteint la bande définie dans l'\cref{ex:m1:exchange-traded-funds:band} (en pointillés) un
\omterm{def:m1:exchange-traded-funds:ap}{participant autorisé} crée ou rachète des parts et la ramène à mi-chemin. À l'intérieur
de la bande, rien ne la ramène à zéro. Données~: simulation du
tutoriel.}\label{fig:m1:exchange-traded-funds:premium}
\end{omfigure}

La plupart des transactions sur ETF n'atteignent jamais le marché primaire. Un \omterm{def:m1:what-a-trading-firm-does:market-maker}{teneur de marché} qui vend
des parts d'ETF se couvre avec des contrats à terme, avec le panier ou avec un ETF corrélé, et
laisse acheteurs et vendeurs se compenser mutuellement~; seul le solde, souvent en fin
de journée, donne lieu à une création. Le nombre de parts en circulation varie peu un
jour où le volume échangé dépasse de plusieurs fois la taille du fonds. Le market making sur les ETF, activité
de plusieurs des sociétés de trading présentées au \cref{ch:m1:what-a-trading-firm-does}, est
traitée dans le livre~11.

\section{Quand la valeur intrinsèque ne peut pas être établie}

L'argument d'arbitrage suppose un cours du panier. Deux types d'ETF n'en disposent
pas. Pour un fonds d'actions japonaises négociées à New York, le marché du panier est
fermé toute la journée. Pour un fonds d'obligations d'entreprises, la plupart des titres
n'ont pas été négociés du tout aujourd'hui~: leurs «~prix~» sont des estimations du fournisseur.

\begin{proposition}[Une valeur liquidative périmée engendre une prime fictive]\label{prop:m1:exchange-traded-funds:stale}
Supposons que la valeur réelle $V_t$ des titres évolue, et que la \omterm{def:m1:exchange-traded-funds:nav}{valeur liquidative} publiée
ne s'ajuste que partiellement chaque jour, $N_t = N_{t-1} + (1-\theta)(V_t - N_{t-1})$ avec
$0 < \theta < 1$. Si l'ETF se négocie à $V_t$, alors après une baisse soutenue de
$V$, l'ETF affiche une décote par rapport à sa \omterm{def:m1:exchange-traded-funds:nav}{valeur liquidative}, qui se résorbe
sans qu'aucune transaction sur l'ETF constitue une bonne affaire.
\end{proposition}
\begin{proof}
$V_t - N_t = \theta(V_t - N_{t-1})$. Après une baisse, $V_t < N_{t-1}$, donc $V_t <
N_t$~: l'ETF, à $V_t$, se situe sous la \omterm{def:m1:exchange-traded-funds:nav}{valeur liquidative}. Lorsque $V$ se stabilise, l'écart
décroît géométriquement au taux $\theta$.
\end{proof}

\begin{example}[Mars 2020]\label{ex:m1:exchange-traded-funds:march}
Lors de la forte baisse de mars 2020, les ETF d'obligations d'entreprises investment grade
ont clôturé environ 365 points de base sous leur \omterm{def:m1:exchange-traded-funds:nav}{valeur liquidative}, en moyenne pondérée par les actifs,
le 12 mars et de nouveau le 19 mars~; le plus gros d'entre eux est passé d'une décote
de 2.8~\% le 20 mars à une \omterm{def:m1:exchange-traded-funds:nav}{prime} de 2.9~\% le 23 mars, jour où une banque centrale
a annoncé un programme d'achats. Une étude de banque centrale consacrée à l'épisode
interprétait les décotes comme le signe que les cours des ETF devançaient des valorisations obligataires périmées, sur un marché
où les dealers avaient cessé de fournir des cours~: l'ETF était le prix, et la
\omterm{def:m1:exchange-traded-funds:nav}{valeur liquidative} n'était qu'une estimation.
\end{example}

\begin{omfigure}
\begin{tikzpicture}
\begin{axis}[width=11cm,height=5cm,
  xlabel={jour},ylabel={valeur par part},
  tick label style={font=\small},label style={font=\small},
  xmin=0,xmax=59,ymin=84,ymax=102,grid=major,grid style={gray!25},
  legend columns=2,legend style={font=\footnotesize,at={(0.5,-0.3)},anchor=north,draw=none,fill=none,/tikz/every even column/.append style={column sep=8pt}},
  table/col sep=comma]
\addplot[omDef,very thick] table[x=day,y=true]{figdata/markets-1/14-exchange-traded-funds/stale.csv};
\addplot[omThm,very thick,dashed] table[x=day,y=nav]{figdata/markets-1/14-exchange-traded-funds/stale.csv};
\legend{cours de l'ETF (valeur réelle des titres),valeur liquidative publiée (valorisations périmées)}
\end{axis}
\end{tikzpicture}

\omcaption{Une baisse simulée sur dix jours dans une classe d'actifs peu liquide. La \omterm{def:m1:exchange-traded-funds:nav}{valeur liquidative}
retarde~; l'ETF semble se négocier avec une décote qui atteint 4~\% puis
«~se résorbe~» à mesure que les valorisations rattrapent leur retard. Données~: simulation du tutoriel, avec
$\theta = 0.75$.}\label{fig:m1:exchange-traded-funds:stale}
\end{omfigure}

La règle pratique~: une \omterm{def:m1:exchange-traded-funds:nav}{prime} n'est une preuve de mauvaise évaluation que dans la mesure où
le panier peut être négocié aux cours utilisés dans la \omterm{def:m1:exchange-traded-funds:nav}{valeur liquidative}. Pour les actions locales,
il le peut, et sept points de base suffisent à tout expliquer. Pour les obligations et
pour les marchés fermés, l'ETF est souvent la meilleure estimation disponible de la valeur, et
la bonne comparaison se fait avec les contrats à terme, avec d'autres ETF et avec les rares obligations
qui ont effectivement été négociées.

\section{Fonds à effet de levier et fonds inverses}

\begin{definition}[ETF à effet de levier]\label{def:m1:exchange-traded-funds:leveraged}
Un \emph{ETF à effet de levier}\index{ETF à effet de levier} de facteur $\beta$ promet
$\beta$ fois le \emph{rendement quotidien} d'un indice, avant frais~; un
\emph{ETF inverse} a $\beta < 0$. Il obtient l'exposition au moyen de swaps et
de contrats à terme, et la réajuste à chaque clôture.
\end{definition}

\begin{proposition}[Le rééquilibrage quotidien]\label{prop:m1:exchange-traded-funds:rebalance}
Un fonds dont l'actif net vaut $N$ et l'exposition $\beta N$ subit un rendement de l'indice
$r$ au cours de la journée. Pour afficher un effet de levier $\beta$ à la clôture, il doit
augmenter son exposition de
\[
\Delta = \beta(\beta - 1)\,N\,r .
\]
Comme $\beta(\beta-1) > 0$ pour tout $\beta$ en dehors de $[0,1]$, les fonds à effet de levier
\emph{et} les fonds inverses achètent tous après une hausse et vendent après une baisse.
\end{proposition}
\begin{proof}
Après le mouvement, l'actif vaut $N(1+\beta r)$ et l'exposition $\beta N(1+r)$.
L'exposition visée est $\beta N(1+\beta r)$. L'écart vaut $\beta N(\beta
r - r)$.
\end{proof}

\begin{example}[Six pour cent de l'actif par pour cent]\label{ex:m1:exchange-traded-funds:six}
Pour $\beta = 3$, $\Delta = 6Nr$ ; pour $\beta = -1$, $\Delta = 2Nr$ ; pour $\beta
= -3$, $\Delta = 12Nr$. Un fonds à effet de levier 3 de 10~milliards de dollars doit acheter
3~milliards de dollars d'exposition à la clôture d'une journée où son indice a progressé de 5~\%, quoi que pense quiconque
au sein du fonds. Le flux relève d'une arithmétique publique, se concentre dans les
dernières minutes de séance et a le signe du mouvement de la journée~: il amplifie les
mouvements de clôture, et anticiper ce flux fait partie des stratégies événementielles du livre~8
(\cref{fig:m1:exchange-traded-funds:rebalance}).
\end{example}

\begin{omfigure}
\begin{tikzpicture}
\begin{axis}[width=9.5cm,height=4.4cm,ybar,bar width=14pt,
  xlabel={facteur d'effet de levier $\beta$},ylabel={transaction (\% de l'actif net)},
  symbolic x coords={-3,-2,-1,1,2,3},xtick=data,
  tick label style={font=\small},label style={font=\small},
  ymin=0,ymax=14,ymajorgrids,grid style={gray!25},
  nodes near coords,nodes near coords style={font=\footnotesize},table/col sep=comma]
\addplot[fill=omDef!60,draw=omDef] table[x=beta,y=trade_pct_nav]{figdata/markets-1/14-exchange-traded-funds/rebalance.csv};
\end{axis}
\end{tikzpicture}

\omcaption{Exposition qu'un fonds doit acheter à la clôture après une journée à $+1\%$, en
pourcentage de son actif net~: $\beta(\beta-1)$. Les fonds inverses échangent davantage que les fonds à effet de levier
de même amplitude, et dans le même sens. Données~:
\cref{prop:m1:exchange-traded-funds:rebalance}.}\label{fig:m1:exchange-traded-funds:rebalance}
\end{omfigure}

\begin{definition}[Érosion par la volatilité]\label{def:m1:exchange-traded-funds:decay}
L'\emph{érosion par la volatilité}\index{erosion@érosion par la volatilité} est l'écart défavorable entre
le rendement sur plusieurs jours d'un fonds à effet de levier et $\beta$ fois (ou la puissance $\beta$-ième de) le
rendement de l'indice sur la même période, causé par la remise à zéro quotidienne
dans un marché volatil.
\end{definition}

\begin{proposition}[Ampleur de l'érosion]\label{prop:m1:exchange-traded-funds:decay}
Si l'indice suit $\dd S/S = \mu\dd t + \sigma\dd W$ et que le fonds est
rééquilibré en continu à l'effet de levier $\beta$, sans coût de financement, alors
\[
\frac{L_T}{L_0} \;=\; \left(\frac{S_T}{S_0}\right)^{\beta}
\exp\!\Bigl(-\tfrac12\,\beta(\beta-1)\,\sigma^2 T\Bigr).
\]
\end{proposition}
\begin{proof}
$\dd L/L = \beta\,\dd S/S$, donc par la formule d'Itô $\dd\ln L = \beta\mu\dd t +
\beta\sigma\dd W - \tfrac12\beta^2\sigma^2\dd t$, tandis que $\beta\,\dd\ln S =
\beta\mu\dd t + \beta\sigma\dd W - \tfrac12\beta\sigma^2\dd t$. On soustrait et
on intègre.
\end{proof}

Avec $\sigma = 30\%$ et $T = 1$ an, le facteur vaut $\exp(-0.27) = 0.76$ pour
$\beta = 3$ et $\exp(-0.09) = 0.91$ pour $\beta = -1$. Le fonds ne triche pas~: il a livré
$\beta$ fois le rendement de chaque journée. C'est l'investisseur
qui le détient pendant un an, en espérant $\beta$ fois le rendement de l'année, qui a
mal lu le contrat.

\begin{omfigure}
\begin{tikzpicture}
\begin{axis}[width=11cm,height=5cm,
  xlabel={jour de séance},ylabel={valeur de 1 investi},
  tick label style={font=\small},label style={font=\small},
  xmin=1,xmax=252,ymin=0.4,ymax=1.8,grid=major,grid style={gray!25},
  legend columns=3,legend style={font=\footnotesize,at={(0.5,-0.3)},anchor=north,draw=none,fill=none,/tikz/every even column/.append style={column sep=8pt}},
  table/col sep=comma]
\addplot[omExo,very thick] table[x=day,y=index]{figdata/markets-1/14-exchange-traded-funds/leveraged.csv};
\addplot[omThm,very thick] table[x=day,y=three_x]{figdata/markets-1/14-exchange-traded-funds/leveraged.csv};
\addplot[omDef,thick,dashed] table[x=day,y=naive_three_x]{figdata/markets-1/14-exchange-traded-funds/leveraged.csv};
\legend{indice,fonds à trois fois la séance,trois fois le rendement de l'indice}
\end{axis}
\end{tikzpicture}

\omcaption{Une année simulée avec une volatilité quotidienne de 2~\% et des rendements quotidiens
dont la somme est nulle. L'indice finit en baisse de 5.8~\%~; trois fois cette baisse serait de
$-17.4\%$~; le fonds à trois fois la séance finit en baisse de 41.8~\%. Données~: la simulation du
tutoriel.}\label{fig:m1:exchange-traded-funds:decay}
\end{omfigure}

\section{Tutoriel~: bandes, valeurs obsolètes et érosion}\label{tut:m1:exchange-traded-funds:bands}

\begin{tutorial}
\textbf{Objectif.} Calculer une bande d'arbitrage, reproduire une décote fallacieuse à partir
d'une \omterm{def:m1:exchange-traded-funds:nav}{valeur liquidative} obsolète, et mesurer l'\omterm{def:m1:exchange-traded-funds:decay}{érosion par la volatilité} par rapport à la formule. \textbf{État
final~:} les quatre figures de données de ce chapitre.
\begin{tutsteps}
\item \textbf{La bande} est une somme de quatre coûts.
\omcode{code/markets-1/14-exchange-traded-funds/python/etf_demo.py}{9}{19}{Les coûts qu'un participant autorisé doit couvrir avant qu'un arbitrage ne soit rentable.}\label{lst:m1:exchange-traded-funds:band}
\item \textbf{\omterm{def:m1:exchange-traded-funds:nav}{Valeur liquidative} obsolète.} Une seule ligne de lissage suffit à fabriquer une décote
de 4~\% dans une baisse de dix jours.
\omcode{code/markets-1/14-exchange-traded-funds/python/etf_demo.py}{39}{45}{Une valeur liquidative calculée à partir de cotations qui ne s'ajustent que partiellement chaque jour.}\label{lst:m1:exchange-traded-funds:stale}
\item \textbf{Le rééquilibrage et l'érosion.}
\omcode{code/markets-1/14-exchange-traded-funds/python/etf_demo.py}{48}{60}{La transaction de clôture d'un fonds à effet de levier, sa trajectoire, et le facteur d'érosion en temps continu.}\label{lst:m1:exchange-traded-funds:lev}
\item \textbf{Vérifier la formule.} Les tests simulent 4\,000 années à une volatilité de 30~\%
et comparent le logarithme moyen du rapport du fonds à trois fois la séance au cube de l'indice
avec $-\tfrac12\cdot6\cdot0.09 = -0.27$.
\end{tutsteps}
\textbf{Pour aller plus loin.} Ajouter au fonds à effet de levier un spread de financement et des frais annuels
de 0.95~\%. Puis faire tendre l'indice (autocorrélation positive des rendements
quotidiens) et trouver le régime dans lequel le fonds à trois fois la séance
\emph{bat} trois fois l'indice.
\end{tutorial}

\section{Construction~: le moniteur de prime}\label{bld:m1:exchange-traded-funds:monitor}

\begin{build}
\bfield{Rôle} La petite société fait du market making sur les ETF de sa bourse simulée
(livre~11). Elle a besoin d'une fair value en temps réel pour chaque fonds, et
d'une alerte lorsque le cours de marché sort de la bande d'arbitrage.
\bfield{Interface} \texttt{Basket(etf, shares\_per\_unit, components, cash)}
avec components sous forme de couples (symbole, quantité par \omterm{def:m1:exchange-traded-funds:ap}{unité de création})~;
\texttt{indicative\_nav(basket, mids)}~; \texttt{premium\_bp(price, inav)}~;
\texttt{band\_bp(costs)}~; \texttt{signal(price, inav, costs)} renvoyant
\texttt{create}, \texttt{redeem} ou \texttt{none}~; et, pour chaque composant, un indicateur
\texttt{stale} lorsque sa dernière cotation est plus ancienne qu'un seuil.
\bfield{Règles} Prix entiers~; la \omterm{def:m1:exchange-traded-funds:nav}{valeur liquidative indicative} utilise les points médians de
la cotation consolidée de la \cref{bld:m1:us-equity-market-structure:nbbo}. Si des composants dont la
valeur dépasse une fraction paramétrable du panier sont obsolètes,
le moniteur renvoie \texttt{unreliable} au lieu d'une \omterm{def:m1:exchange-traded-funds:nav}{prime}. Les opérations sur titres
modifient le panier pendant la nuit
(\cref{bld:m1:shares-corporate-actions-indices:adjuster}).
\bfield{Tests d'acceptation} \texttt{code/firm/etfmonitor/tests/}~: un panier de trois actions évalué
à la main~; des signaux à la bande et au-delà~; la règle des composants
obsolètes.
\bfield{Pour aller plus loin} Remplacer les composants dont le marché est fermé par une régression sur
des instruments ouverts (contrats à terme, devise, un ETF apparenté)~: la fair value
d'un fonds international.
\end{build}

\begin{omsources}
\item ETFGI, communiqués de presse sur les actifs de l'industrie mondiale des ETF, juillet--septembre
2026 (selon Markets Media, «~Global ETF Assets Top \$24 Trillion for
First Time~»).
\item US Securities and Exchange Commission, \emph{SEC Adopts New Rule to
Modernize Regulation of Exchange-Traded Funds}, communiqué de presse 2019-190, septembre 2019.
\item S.~Aramonte et F.~Avalos, «~The recent distress in corporate bond
markets: cues from ETFs~», \emph{BIS Bulletin} n°~6, 14 avril 2020; K.~Todorov,
«~The anatomy of bond ETF arbitrage~», \emph{BIS Quarterly Review}, mars
2021.
\item M.~Cheng et A.~Madhavan, «~The dynamics of leveraged and inverse
exchange-traded funds~», \emph{Journal of Investment Management} 7 (2009).
\item A.~Madhavan, \emph{Exchange-Traded Funds and the New Dynamics of
Investing}, Oxford University Press, 2016.
\end{omsources}

\section{Exercices}

\begin{exercise}[$\star$]\label{exo:m1:exchange-traded-funds:1}
La \omterm{def:m1:exchange-traded-funds:nav}{valeur liquidative} d'un ETF est de 82.40~\$ et celui-ci se négocie à 82.31~\$. Donner la \omterm{def:m1:exchange-traded-funds:nav}{prime} en
points de base. Avec une bande de 9 points de base, un \omterm{def:m1:exchange-traded-funds:ap}{participant autorisé} intervient-il,
et comment~?
\end{exercise}

\begin{exercise}[$\star$]\label{exo:m1:exchange-traded-funds:2}
Une \omterm{def:m1:exchange-traded-funds:ap}{unité de création} représente 50\,000 actions d'ETF à une \omterm{def:m1:exchange-traded-funds:nav}{valeur liquidative} de 40~\$, et le fonds facture
500~\$ par création. Exprimer les frais en points de base de la valeur de l'unité. Quel
est ce montant si le \omterm{def:m1:exchange-traded-funds:ap}{participant autorisé} crée dix unités en un seul ordre pour les mêmes frais~?
\end{exercise}

\begin{exercise}[$\star$]\label{exo:m1:exchange-traded-funds:3}
Un fonds $-2\times$ dispose d'un actif net de 800~millions de dollars et son indice baisse de 4~\%
aujourd'hui. Donner le rendement du fonds, son nouvel actif net et la transaction
qu'il doit réaliser à la clôture, en précisant son sens.
\end{exercise}

\begin{exercise}[$\star\star$]\label{exo:m1:exchange-traded-funds:4}
Un indice progresse de 10~\% le premier jour, puis recule de 9.09~\% le second, et termine
à son niveau de départ. Calculer le rendement sur deux jours d'un fonds à $2\times$, d'un fonds à $3\times$ et
d'un fonds à $-1\times$. Lequel perd le plus, et pourquoi~?
\end{exercise}

\begin{exercise}[$\star\star$]\label{exo:m1:exchange-traded-funds:5}
Pour $\sigma = 40\%$ par an, calculer le facteur d'érosion de la
\cref{prop:m1:exchange-traded-funds:decay} sur un an pour $\beta = 2$, $3$
et $-2$. Sur quel horizon un fonds à $3\times$ perd-il 10~\% par érosion à cette
volatilité~?
\end{exercise}

\begin{exercise}[$\star\star$]\label{exo:m1:exchange-traded-funds:6}
La \omterm{def:m1:exchange-traded-funds:nav}{valeur liquidative} d'un ETF obligataire suit la \cref{prop:m1:exchange-traded-funds:stale} avec
$\theta = 0.8$. La valeur vraie tombe de 100 à 94 en un jour et s'y maintient. Donner la
\omterm{def:m1:exchange-traded-funds:nav}{valeur liquidative} et la décote apparente ce jour-là et
les trois jours suivants. Après combien de jours la décote est-elle inférieure à 0.5~\%~?
\end{exercise}

\begin{exercise}[$\star\star\star$]\label{exo:m1:exchange-traded-funds:7}
\emph{Programmation.} Avec \texttt{simulate\_premium} et les coûts définis
à l'\cref{ex:m1:exchange-traded-funds:band}, lancer 5\,000 pas avec la graine 1 et
indiquer le nombre de pas auxquels un \omterm{def:m1:exchange-traded-funds:ap}{participant autorisé} est intervenu et la
\omterm{def:m1:exchange-traded-funds:nav}{prime} absolue moyenne. Diviser par deux le coût du panier $c_B$ et indiquer la même grandeur. Que
gagne un investisseur dans l'ETF~?
\end{exercise}

\begin{exercise}[$\star\star\star$]\label{exo:m1:exchange-traded-funds:8}
\emph{Trouver la faille.} Une lettre d'information écrit~: «~Cet ETF obligataire à haut rendement a clôturé
avec une décote de 3~\% par rapport à sa \omterm{def:m1:exchange-traded-funds:nav}{valeur liquidative}, sur un marché en repli. Il faut l'acheter~: la décote
se résorbe toujours en deux semaines, un gain sans risque de 3~\%.~» Expliquer pourquoi la résorption
de la décote n'implique pas de profit, et décrire la position qui serait nécessaire pour
la capturer si elle était réelle.
\end{exercise}

\section{Problème~: le rééquilibrage quotidien d'un fonds à triple effet de levier}

\begin{problem}[{Problème du week-end --- les dix dernières minutes d'un fonds à effet de levier}]\label{pb:m1:exchange-traded-funds:1}
Un fonds $3\times$ et un fonds $-3\times$ répliquent le même indice d'actions. À la clôture
d'hier soir, le premier avait un actif net de 8~milliards de dollars et le second de 3
milliards de dollars. Les contrats à terme sur indice s'échangent pour 300~milliards de dollars par jour, dont un dixième dans les
dix dernières minutes.

\medskip\noindent\textbf{Partie~I --- Une journée calme.}
\begin{enumerate}
\item L'indice progresse de 1~\%. Donner le rendement et le nouvel actif net de chaque fonds.
\item Donner la transaction de clôture de chaque fonds avec son sens, ainsi que le total.
\item Exprimer le total en fraction du volume des dix dernières minutes.
\item Avec la règle de la racine carrée de la \cref{met:m1:the-sell-side:sqrt}, une volatilité quotidienne
de 1~\% et $\kappa = 0.7$, estimer l'impact de cette transaction,
en prenant $V$ pour le volume des dix minutes et en ramenant $\sigma$ à dix minutes
d'une journée de 390 minutes.
\end{enumerate}

\medskip\noindent\textbf{Partie~II --- Une journée violente.}
\begin{enumerate}[resume]
\item L'indice recule de 5~\%. Donner le rendement et le nouvel actif net de chaque fonds.
\item Donner les transactions de clôture et le total.
\item En fraction du volume des dix dernières minutes~?
\item Estimer l'impact comme à la question 4 avec une volatilité quotidienne de 3~\%.
\item L'impact se répercute en retour~: si les ventes des fonds font encore bouger l'indice
de $x$, elles doivent vendre un peu plus. En prenant la réponse à la question 8 comme $x$,
de combien la transaction requise augmente-t-elle~?
\end{enumerate}

\medskip\noindent\textbf{Partie~III --- Une mauvaise semaine.}
L'indice présente des rendements de $-5\%$, $+4\%$, $-6\%$, $+5\%$, $+2.5\%$ sur cinq jours.
\begin{enumerate}[resume]
\item Donner le rendement cumulé de l'indice.
\item Donner le rendement cumulé du fonds $3\times$, et trois fois celui de
l'indice.
\item Donner le rendement cumulé du fonds $-3\times$, et moins trois fois celui de
l'indice.
\item Un investisseur détient les deux fonds pour des montants égaux en dollars, «~pour être couvert~».
Donner son rendement.
\item Expliquer son résultat à l'aide de la \cref{prop:m1:exchange-traded-funds:decay}.
\end{enumerate}

\medskip\noindent\textbf{Partie IV --- Jugement.}
\begin{enumerate}[resume]
\item Qui se trouve de l'autre côté des transactions de clôture de la partie~II, et
que doivent-ils savoir à l'avance~?
\item Pourquoi le fonds $-3\times$ échange-t-il deux fois plus par dollar d'actifs
que le fonds $3\times$~?
\item Quel mouvement quotidien de l'indice effacerait le fonds $3\times$~? Comment les
prospectus de ces fonds traitent-ils ce cas~?
\item Un régulateur demande si ces fonds déstabilisent la clôture. Donner la grandeur qui
répond à la question et sa valeur le jour de forte baisse.
\item Énoncer le \emph{résultat nommé}~: le notionnel que les deux fonds doivent échanger
ensemble à la clôture après la baisse de 5~\%, en milliards de dollars.
\item En une phrase~: pour qui un \omterm{def:m1:exchange-traded-funds:leveraged}{ETF à effet de levier} est-il un instrument judicieux~?
\end{enumerate}
\end{problem}

\section{Questions d'entretien}

\begin{interviewq}[$\star$ \iqroles{trader, chercheur, développeur}]\label{iq:m1:exchange-traded-funds:1}
Comment le cours d'un ETF reste-t-il proche de la valeur de ce qu'il détient~?
\end{interviewq}

\begin{interviewq}[$\star$ \iqroles{trader, chercheur}]\label{iq:m1:exchange-traded-funds:2}
L'indice d'un \omterm{def:m1:exchange-traded-funds:leveraged}{ETF à effet de levier} $2\times$ progresse de 10~\% sur l'année. L'ETF progresse-t-il
de 20~\%~?
\end{interviewq}

\begin{interviewq}[$\star\star$ \iqroles{trader, chercheur}]\label{iq:m1:exchange-traded-funds:3}
Un ETF détenant des actions japonaises se négocie à New York avec une \omterm{def:m1:exchange-traded-funds:nav}{prime} de 1.5~\% sur sa
\omterm{def:m1:exchange-traded-funds:nav}{valeur liquidative} à 11:00, heure de New York. S'agit-il d'un arbitrage~?
\end{interviewq}

\begin{interviewq}[$\star\star$ \iqroles{trader, chercheur}]\label{iq:m1:exchange-traded-funds:4}
Vous êtes \omterm{def:m1:what-a-trading-firm-does:market-maker}{teneur de marché} sur un ETF de 500 actions. Un client vous achète pour
50~millions de dollars. Décrivez-moi votre couverture et votre fin de journée.
\end{interviewq}

\begin{interviewq}[$\star\star$ \iqroles{chercheur, ingénieur ML}]\label{iq:m1:exchange-traded-funds:5}
Établissez le montant qu'un fonds à effet de levier doit échanger à la clôture en fonction
de son levier et du rendement du jour. Qu'y a-t-il de surprenant dans les fonds inverses~?
\end{interviewq}

\begin{interviewq}[$\star\star\star$ \iqroles{chercheur, trader}]\label{iq:m1:exchange-traded-funds:6}
En pleine crise, un ETF obligataire se négocie 4~\% en dessous de sa \omterm{def:m1:exchange-traded-funds:nav}{valeur liquidative}. Construisez l'argument selon lequel l'ETF est
mal évalué, celui selon lequel la \omterm{def:m1:exchange-traded-funds:nav}{valeur liquidative} l'est, et dites quelles données trancheraient.
\end{interviewq}
